\section{A universal phase-sector lower selection theorem}
\label{sec:sector-theorem}

We first remove all arithmetic.  Let \(\cU\) be a finite set partitioned
into fibers
\begin{equation}
 \cU=\bigsqcup_{y\in\cY}\cF_y,
 \label{eq:tpc56-abstract-fibers}
\end{equation}
and let \(x=(x_u)_{u\in\cU}\in\C^{\cU}\).  The unweighted grouping map is
\begin{equation}
 (\Sigma x)_y=\sum_{u\in\cF_y}x_u.
 \label{eq:tpc56-abstract-grouping}
\end{equation}
Choose an arbitrary unit gauge \(\omega_y\in\C\),
\(|\omega_y|=1\), on every fiber.  The gauge family is part of the data and
is fixed before the maximizing sector index is selected.  These gauges only
rotate output coordinates and hence do not change any grouped norm.

\subsection{Equal sectors and one global choice}

Fix an integer \(M\ge3\).  On the phase circle use the half-open sectors
\begin{equation}
 I_k=
 \left[\frac{2\pi k}{M}-\frac\pi M,
       \frac{2\pi k}{M}+\frac\pi M\right)
 \pmod {2\pi},
 \qquad 0\le k<M.
 \label{eq:tpc56-equal-sectors}
\end{equation}
For \(x_u\ne0\), put
\begin{equation}
 \varphi_u=\Arg(\overline{\omega_{q(u)}}x_u)\pmod {2\pi},
 \label{eq:tpc56-normalized-phase}
\end{equation}
where \(q(u)=y\) for \(u\in\cF_y\), and define
\begin{equation}
 \cS_k(x;\omega)=\{u:x_u\ne0,\ \varphi_u\in I_k\}.
 \label{eq:tpc56-sector-support}
\end{equation}
Let \(P_k=P_k(x;\omega)\) be the orthogonal coordinate projection onto
\(\cS_k(x;\omega)\).  Zero coordinates can be assigned to any sector;
they affect none of the identities.  The notation records an important
fact: after \(x\) is fixed, \(P_k\) is linear and orthogonal, but the rule
\(x\mapsto P_k(x;\omega)\) is nonlinear.

Define the selected diagonal and grouped energies
\begin{equation}
 D_k(x):=\|P_kx\|_2^2,
 \qquad
 C_k(x):=\|\Sigma P_kx\|_2^2.
 \label{eq:tpc56-sector-energies}
\end{equation}

\begin{theorem}[Equal-sector compression and selection]
\label{thm:tpc56-equal-sector-selection}
For every finite fiber system, every \(x\), every fixed family of unit
gauges, and every integer \(M\ge3\),
\begin{align}
 \sum_{k=0}^{M-1}D_k(x)
 &=\|x\|_2^2,
 \label{eq:tpc56-sector-diagonal-partition}\\
 C_k(x)&\ge\cos^2(\pi/M)D_k(x)
 \qquad(0\le k<M),
 \label{eq:tpc56-each-sector-lower}\\
 \sum_{k=0}^{M-1}C_k(x)
 &\ge\cos^2(\pi/M)\|x\|_2^2.
 \label{eq:tpc56-sector-compression-lower}
\end{align}
Let \(k_*\) be the least index maximizing \(D_k(x)\).  Then this canonical
tie rule gives
\begin{equation}
 \boxed{
 \|\Sigma P_{k_*}x\|_2^2
 \ge c_M\|x\|_2^2,
 \qquad
 c_M:=\frac{\cos^2(\pi/M)}M.}
 \label{eq:tpc56-universal-sector-lower}
\end{equation}
The same index \(k_*\) is used on every fiber in the normalized coordinates
\(\Arg(\overline{\omega_y}x_u)\).  Since \(\omega_y\) may vary with \(y\),
the corresponding absolute angular window is a fiber-dependent rotation.
\end{theorem}

\begin{proof}
The sector supports partition the nonzero coordinates, proving
\eqref{eq:tpc56-sector-diagonal-partition}.  Fix \(k\) and a fiber \(y\).
For every \(u\in\cF_y\cap\cS_k\), one can write
\begin{equation}
 \overline{\omega_y}e^{-2\pi i k/M}x_u
 =|x_u|e^{i\delta_u},
 \qquad |\delta_u|\le\frac\pi M,
 \label{eq:tpc56-sector-angle-form}
\end{equation}
with one endpoint understood by the half-open convention.  Therefore
\begin{align}
 \left|\sum_{u\in\cF_y\cap\cS_k}x_u\right|
 &\ge
 \operatorname{Re}\left(
 \overline{\omega_y}e^{-2\pi i k/M}
 \sum_{u\in\cF_y\cap\cS_k}x_u\right)\notag\\
 &\ge\cos(\pi/M)
 \sum_{u\in\cF_y\cap\cS_k}|x_u|\notag\\
 &\ge\cos(\pi/M)
 \left(\sum_{u\in\cF_y\cap\cS_k}|x_u|^2\right)^{1/2}.
 \label{eq:tpc56-fiber-sector-proof}
\end{align}
Squaring and summing over \(y\) proves
\eqref{eq:tpc56-each-sector-lower}; summing over \(k\) gives
\eqref{eq:tpc56-sector-compression-lower}.  Finally,
\(D_{k_*}\ge M^{-1}\|x\|_2^2\), so
\eqref{eq:tpc56-each-sector-lower} gives
\eqref{eq:tpc56-universal-sector-lower}.
\end{proof}

The theorem contains two different objects.  The phase-split synthesis
\begin{equation}
 x\longmapsto
 \left(\sum_{u\in\cF_y\cap\cS_k}x_u\right)_{(y,k)}
 \label{eq:tpc56-phase-split-synthesis}
\end{equation}
has a lower bound on the distinguished vector used to define its sectors.
The original output instead sums the \(M\) sector outputs at fixed \(y\).
Those final \(M\) channels may cancel.  Thus
\eqref{eq:tpc56-sector-compression-lower} is not a lower bound for
\(\|\Sigma x\|_2\).

\subsection{Why five equal sectors are best in this family}

\begin{proposition}[Optimal equal-sector count]
\label{prop:tpc56-optimal-equal-sector-count}
Among integers \(M\ge3\), the certified constant
\(c_M=M^{-1}\cos^2(\pi/M)\) is uniquely maximized at \(M=5\).  Its value is
\begin{equation}
 \boxed{
 c_5=\frac{\cos^2(\pi/5)}5
     =\frac{3+\sqrt5}{40}
     =0.1309016994\ldots .}
 \label{eq:tpc56-pentant-constant}
\end{equation}
\end{proposition}

\begin{proof}
For real \(t\ge3\), let
\(f(t)=t^{-1}\cos^2(\pi/t)\).  Direct differentiation gives
\begin{equation}
 \frac{f'(t)}{f(t)}
 =\frac{2\pi\tan(\pi/t)-t}{t^2}.
 \label{eq:tpc56-equal-sector-derivative}
\end{equation}
After setting \(\theta=\pi/t\), the unique critical point is determined by
\(2\theta\tan\theta=1\).  The left side is strictly increasing on
\((0,\pi/2)\), and the corresponding \(t\) lies strictly between \(4\)
and \(5\).  Hence the integer maximum is either \(4\) or \(5\).  Now
\begin{equation}
 c_4=\frac18,
 \qquad
 c_5=\frac{3+\sqrt5}{40}>\frac18.
 \label{eq:tpc56-four-five-comparison}
\end{equation}
The monotonicity on either side of the critical point proves uniqueness.
\end{proof}

The word \emph{optimal} in \cref{prop:tpc56-optimal-equal-sector-count}
refers only to the constants certified by equal finite sectorizations.  It
is not a global optimality theorem over all nonlinear subset rules.

\subsection{A continuously rotated window}

The finite pentant library is convenient for a literal audit.  If a
continuum of rotations is allowed, the constant can be improved slightly.
For \(0<\theta<\pi/2\) and \(\alpha\in\T=\R/(2\pi\mathbb Z)\), retain the
coordinates satisfying
\begin{equation}
 \dist_{\T}(\varphi_u,\alpha)\le\theta
 \label{eq:tpc56-rotating-window}
\end{equation}
and denote the associated coordinate projection by \(P_{\alpha,\theta}\).

\begin{theorem}[Rotating-window selection]
\label{thm:tpc56-rotating-window}
For every \(0<\theta<\pi/2\), there exists \(\alpha\in\T\) such that
\begin{equation}
 \|\Sigma P_{\alpha,\theta}x\|_2^2
 \ge c(\theta)\|x\|_2^2,
 \qquad
 c(\theta):=\frac\theta\pi\cos^2\theta.
 \label{eq:tpc56-rotating-window-bound}
\end{equation}
The function \(c(\theta)\) has a unique maximizer
\(\theta_*\in(0,\pi/2)\) characterized by
\begin{equation}
 2\theta_*\tan\theta_*=1.
 \label{eq:tpc56-optimal-window-equation}
\end{equation}
Numerically,
\begin{equation}
 \theta_*=0.6532711871\ldots,
 \qquad
 c_*=c(\theta_*)=0.1311275284\ldots .
 \label{eq:tpc56-optimal-window-constant}
\end{equation}
\end{theorem}

\begin{proof}
For each nonzero coordinate, the set of centers \(\alpha\) satisfying
\eqref{eq:tpc56-rotating-window} has normalized Haar measure
\(\theta/\pi\).  Fubini therefore gives
\begin{equation}
 \int_{\T}\|P_{\alpha,\theta}x\|_2^2
 \frac{d\alpha}{2\pi}
 =\frac\theta\pi\|x\|_2^2.
 \label{eq:tpc56-window-average-energy}
\end{equation}
Some center retains at least that much diagonal energy.  The proof of
\eqref{eq:tpc56-fiber-sector-proof}, with \(\theta\) in place of
\(\pi/M\), proves \eqref{eq:tpc56-rotating-window-bound}.  Finally,
\begin{equation}
 c'(\theta)=\frac{\cos\theta}{\pi}
 \bigl(\cos\theta-2\theta\sin\theta\bigr),
 \label{eq:tpc56-window-derivative}
\end{equation}
so the unique interior critical point satisfies
\eqref{eq:tpc56-optimal-window-equation}; endpoint values are zero.
\end{proof}

\subsection{Stability under phase uncertainty}

The sector theorem is not destroyed by small errors in the phases used to
form the coordinate sets.

\begin{theorem}[Phase-error stability]
\label{thm:tpc56-phase-error-stability}
Suppose approximate normalized phases \(\widehat\varphi_u\) satisfy
\begin{equation}
 \dist_{\T}(\widehat\varphi_u,\varphi_u)\le\varepsilon
 \qquad(x_u\ne0).
 \label{eq:tpc56-phase-error-hypothesis}
\end{equation}
Partition the coordinates by the \(M\) equal sectors of the approximate
phases, and choose the least sector with maximal true diagonal energy.
If
\begin{equation}
 0\le\varepsilon<\frac\pi2-\frac\pi M,
 \label{eq:tpc56-phase-error-range}
\end{equation}
then the resulting projection \(\widehat P_*\) satisfies
\begin{equation}
 \boxed{
 \|\Sigma\widehat P_*x\|_2^2
 \ge
 \frac{\cos^2(\pi/M+\varepsilon)}M\|x\|_2^2.}
 \label{eq:tpc56-phase-error-lower}
\end{equation}
\end{theorem}

\begin{proof}
The approximate sectors still partition the active coordinates, so the
selected sector carries at least \(M^{-1}\|x\|_2^2\) diagonal energy.
Every true phase in that sector lies within
\(\pi/M+\varepsilon<\pi/2\) of its center.  Apply the real-projection
argument in \eqref{eq:tpc56-fiber-sector-proof} with this enlarged
half-angle.
\end{proof}

The error \(\varepsilon\) may include both an atomwise phase error and a
fiberwise error in the chosen common gauge.  The theorem requires exact
magnitudes only to identify an energy-maximizing sector; a separate
magnitude-approximation error would have to be charged explicitly.
